\lecture{13}{Fri 27 Mar 2026 14:01}{Free field theories, continued}

Let us recall what is happening. Earlier in the course, we took a na\"ive approach to free field theories. We started with a vector bundle $E\to M$ on a Riemannian manifold $M$. The \emph{fields} were then defined as the sheaf of sections $\mathcal{E} $ of $E$. The action functional was then defined as a map $S : \mathcal{E}  \to \mathbb{R} $. Here we do not claim that $S$ is a map of (pre)sheaves, and simply use this notation to indicate that for all open $U\subseteq M$, there is a map $S : \mathcal{E} (U) \to \mathbb{R} $. If one passes to compactly supported sections, then $S$ is actually a map of cosheaves $S : \mathcal{E} _c \to \mathbb{R} $.

Physically, the equations of motion are extrema for the action functional. Thus we define the \emph{Euler-Lagrange equations of motion} as the critical locus $\operatorname{Crit}(S)=\mathrm{d} S^{-1} (0) $. Here we recall that $\mathrm{d} S : T\mathcal{E}(U) \to \mathbb{R} $. Here we have viewed $\mathcal{E} (U)$ as a vector space, and thus a smooth manifold.

From here, there are two perspectives we can take: the classical perspective, and the quantum perspective. In the quantum perspective, we want to define correlation functions by integrating against the measure $e^{S/\hbar } \omega _0$, where $\omega _0$ is a volume form. The quantum perspective is highly involved, so we continue for now with the classical perspective to build the necessary background.

In the classical perspective, observable quantities are functions $\operatorname{Crit} (S) \to \mathbb{R} $. We often work with polynomials for a important yet highly explicit example. Polynomials are informally understood to be defined as elements of the symmetric algebra of the dual space $\operatorname{Sym} (\operatorname{Crit} (S)^*)$.

Before we continue, there is a simplification we can provide. A dg manifold $M$ is, intuitively, a manifold $M$ with a nice sheaf of dg commutative algebras. For example, the derived critical locus $(M,\mathrm{PV} ^{\bullet} )$ is a dg manifold. Here, ``nice'' means it is sufficiently close to a symmetric algebra.

We can establish some criterion for when such a dg manifold is \emph{linear}. It is linear if $M=V$ is a vector space, and if the action is a quadratic form. In our case, the base manifold is the space of fields $C^\infty (M)$, which is a vector space, and the differential is indeed a quadratic form.

If $M$ is a vector space, then $TM$ is trivial, so sections are just maps $M\to M$. We note that $C^\infty (M,M)$ can be identified with $C^\infty (M)\otimes M$ by taking a smooth map $f : M \to M$, choosing a basis $\left\{ e_i \right\} _{1\leq i\leq n} $, and writing $f=\sum_{i} f_ie_i$. Then simply map $f\mapsto \sum_{i} f_i\otimes e_i$.

More generally, triviality of the tangent bundle yields that $k$-polyvector fields are maps $M\to M\times \Lambda ^kM$, which are similarly identified with $C^\infty (M)\otimes \Lambda ^kM$. We then use $\Lambda ^{\bullet} M \cong \operatorname{Sym} (M[1])$ to obtain
\[
	\mathrm{PV} ^{\bullet} \cong C^\infty (M)\otimes \operatorname{Sym} (M[1])
.\]
We assume the situation is sufficiently nice such that $C^\infty (M) \cong \operatorname{Sym} (M^*)$. Then distributing $\operatorname{Sym} $ over tensor yields
\[
	\mathrm{PV}^{\bullet} \cong \operatorname{Sym} (M^* \oplus M[1])
.\]
This is generated by the chain complex $M^* \oplus M[1]$, and we can take the graded dual $(M^*\oplus M[1])^{\vee } $ to obtain a cochain complex $M\oplus M^*[-1]$. We can go backwards, taking the symmetric algebra of the graded dual of any cochain complex, to obtain a linear dg manifold. Thus we can simplify the presentation of a linear dg manifold into a single cochain complex. We call such a complex the \emph{derived space of solutions to the Euler-Lagrange equations}.

\begin{discussion}\label{4.1.3}
	Why do we call this complex the derived space of solutions to the Euler-Lagrange equations? First, we note that the graded vector space $M\oplus M^*[-1]$ without the structure of the cochain complex is referred to as the \emph{derived space of fields}. Recall that we defined $\mathcal{O} (\operatorname{Crit} (S))$ as the Koszul complex $\mathrm{PV} ^{\bullet} \cong C^\infty (M) \otimes \Lambda^{\bullet} V$. The derived critical locus is really the ``functions'' on the critical locus. However, it is unclear how to interpret the space $\mathrm{PV} ^{\bullet} $ as functions on anything, so we define a derived notion of the underlying space $\operatorname{Crit} (S)$. Of course the term ``derived critical locus'' is already taken, so we define such a space to be the \emph{derived space of solutions}.

	The idea is that functions on such a derived space should be equivalent to the derived critical locus, which was defined to somewhat be functions on $\operatorname{Crit} (S)$. We intuitively understand functions on the space as the symmetric algebra of the graded dual, so we simply want a complex $W$ whose symmetric algebra on its graded dual is the derived critical locus $PV^{\bullet} $. As we just saw, in the case where $M=V$ is a vector space, this is solved by $M \to M^*[-1]$. The differential $d : M \to M^*[-1]$ is given by $d(v)=\frac{1}{2} S(v,-)$, where $S(-,-)$ is the bilinear form determined by a quadratic action $S$ on a vector space. This is explained in the following example.

	So why do we call the underlying graded vector space of the derived space of solutions the \emph{derived space of fields}? It will become more clear when we consider the example of abelian Yang-Mills, but for now we are content to simply say that the cohomology $H^0(W)$ of the derived space of solutions is the classical critical locus $\operatorname{Crit} (S)$. Hence adding in the differential corresponds in some way to working with the action functional, and the action functional acts on fields. Hence it makes sense to call the underlying graded vector space the \emph{derived space of fields}.
\end{discussion}

\begin{example}\label{4.1.4}
Let's consider a toy example. Let $V$ be a vector space, and $q : V \to \mathbb{R} $ a quadratic function on $V$. Let $\omega _0$ be the Lebesgue measure on $V$. We want to understand the divergence complex for the meaure $e^{q/\hbar } \omega _0$. We do not need to assume $q$ is nondegenerate here.

We claim that the derived critical locus is described by the derived space of solutions, which we refer to as $W$:
\[\begin{tikzcd}[row sep = 2.5em, column sep = 2.5em]
	V\ar[" \partial / \partial q ", r]  & V^*[-1]
\end{tikzcd}\]
which maps $v\in V$ to the linear functional $\frac{\partial }{\partial q} v=q(v,-)$. Here we abuse notation and identify $q$ with its \emph{polarization identity}. A quadratic form $q$ uniquely determines a bilinear form $B$ defined by $B(v,w) = q(v+w) - q(v) - q(w)$. We abuse notation and write $B=q$.

Note that $W$ comes with a graded antisymmetric pairing $\left\langle -,- \right\rangle : W\otimes W \to \mathbb{R} $ of cohomological degree $-1$. Note that the fact that the pairing has cohomological degree $-1$ means that $\left\langle x,y \right\rangle $ only lands in $\mathbb{R} $, which is concentrated in degree 0, if $\left| x \right| +\left| y \right| =1$. Hence we take $x=v\in V$ and $y=\alpha \in V^*[-1]$, and define
\[
	\left\langle v,\alpha \right\rangle =\alpha (v) = -\left\langle \alpha ,v \right\rangle 
.\]

Now define the (shifted) \emph{Heisenberg} dg Lie algebra $\mathcal{H} _W$ as the dg Lie algebra
\[
	\mathcal{H} _W\coloneqq \mathbb{C} \cdot \hbar [-1]\oplus W,
\]
where $\mathbb{C} \cdot \hbar $ is the 1-dimensional vector space with basis $\hbar $. The Lie bracket is given by
\[
	[w,w']=\hbar \left\langle w,w' \right\rangle ,
\]
and it vanishes on $\mathbb{C} \cdot \hbar $.

Consider the Chevalley-Eilenberg chains on this dg Lie algebra $C_{\bullet} (\mathcal{H} _W)$. Recall that this is defined to be the symmetric algebra of the underlying graded vector space $\mathcal{H} _W[1]=W[1]\oplus \mathbb{C} \cdot \hbar $, whose differential is given by the Lie bracket acting as a coderivation. Note that
\[
	W^\vee =(V\oplus V^*[-1])^\vee =V^*\oplus V[1] = W[1]
\]
We then obtain
\[
	C_{\bullet} (\mathcal{H} _W) \cong \operatorname{Sym} (W^\vee \oplus (\mathbb{C} \cdot \hbar )) \cong \operatorname{Sym} (W^\vee )\otimes \operatorname{Sym} (\mathbb{C} \cdot \hbar )
.\]
Since $\mathbb{C} \cdot \hbar $ is concentrated in degree 0, it symmetric algebra is just the polynomial ring $\mathbb{C} [\hbar ]$, so tensoring with the polynomial ring just adjoins the variable $\hbar $ to the symmetric algebra, yielding
\[
	C_{\bullet} (\mathcal{H} _W) \cong \operatorname{Sym} (W^\vee )[\hbar ]
.\]
We can show that $C_{\bullet} (\mathcal{H} _W)\cong \mathrm{PV} ^{\bullet} (V)[\hbar ]$ by showing that $\operatorname{Sym} (W^\vee )\cong \mathrm{PV} ^{\bullet} (V)$. We have
\[
	\operatorname{Sym} (W^\vee )= \operatorname{Sym} (V^*\oplus V[1]) \cong \operatorname{Sym} (V^*)\otimes \operatorname{Sym} (V[1])\cong \operatorname{Sym} (V^*)\otimes \Lambda ^{\bullet} V,
\]
where we have noted of course that since $V^*$ is concentrated in degree 0, so is its symmetric algebra. Note obviously that $\operatorname{Sym} (V^*)$ consists of polynomials,  and thus $\operatorname{Sym} (V^*)\otimes \Lambda ^{\bullet} V$ can be identified with polyvector fields on $V$.

We see that if we can find the proper derived space of solutions, we should be able to obtain a graded antisymmetric pairing that induces a Heisenberg (dg) Lie algebra whose Chevalley-Eilenberg chains are the derived critical locus.
\end{example}



\section{The prefactorization algebra of a free field theory}

\begin{example}\label{4.2.1}	
Let's consider the more complex example of a basic scalar field theory, which takes place on a Riemannian manifold $(M,g)$, has a space of fields $C^\infty (M)$, and action functional
\[
	S(\phi )=\int_{M} \phi \Delta \phi ,
\]
where $\Delta $ is the Laplacian. We integrate over $U\subseteq M$ when we want to define more local fields. This is much more complex since polyvector fields no longer look like $\Lambda ^{\bullet} M$, but rather are sections of $\Lambda ^{\bullet} TM$. \Cref{4.1.4} suggests that we should be able to define the derived critical locus as the Chevalley-Eilenberg chains of a Heisenberg algebra associated to a cochain complex with a graded antisymmetric pairing; in particular we should find a derived space of solutions to the equations of motion. Consider the sheaf $\mathcal{E} $ valued in cochain complexes
\[
	\mathcal{E}(U) = \left( 
		\begin{tikzcd}
			0\ar[r] & C^\infty (U)\ar[" \Delta  ", r] & C^\infty (U)[-1] \ar[r] & 0
		\end{tikzcd}
	\right) ,
\]
where the leftmost $C^\infty (U)$ is in degree 0 and the rightmost $C^\infty (U)$ is in degree 1. In the future, if $V$ is a vector space, we will write $V[-n]$ to denote that $V$ is in degree $n$. This notation is reconciled with the usual shift functor and the convention of taking vector spaces to be concentrated in degree 0 by simply taking cochain complexes to be valued in cochain complexes. We will also omit the trailing zeroes, and leave them implicit.

Earlier, our derived space of solutions was only one cochain complex. We could extend it to a sheaf of cochain complexes by restricting along open subsets, and the derived critical locus would then be the Chevalley-Eilenberg chains associated to its global sections. However, the derived critical locus here really means the global sections of the derived critical locus, since it is itself a sheaf. Hence we should be able to recover the derived critical locus sheaf by applying this construction pointwise. We proceed by considering the global sections, but there is nothing that prevents us from considering an open subset $U\subseteq M$ instead. Hence, we would like to define observables as the symmetric algebra of the graded dual of $\mathcal{E} $, possibly up to some shift. We can make this idea precise.

We work in a bit more generality than required for this particular problem. Given a graded vector bundle $E = \left\{ E_i \to M \right\}_{i\in\mathbb{Z} } $ over a manifold $M$, define the graded vector bundle $E^!$ on $M$ by taking the (graded) tensor product with the top exterior bundle $E^! \coloneqq E^\vee \otimes_{\mathbb{R}} \Lambda ^nT^*M$, where $\Lambda ^nT^*M$ is concentrated in degree 0. We denote the sheaf of sections of $E$ by $\mathcal{E} $, and the sheaf of sections of $E^!$ by $\mathcal{E} ^!$. We note a theorem: the functor mapping vector bundles over $M$ to their sheaves of sections preserves limits and colimits, and moreover one does not have to sheafify the colimits of presheaves. This holds in the graded case as well, where our sheaves now take values in graded vector spaces. This gives $\mathcal{E} ^! \cong \mathcal{E} ^{\vee } \otimes_{C^\infty} \Omega _M^n$. We denote by $\mathcal{E} _c$ the cosheaf of compactly supported sections of $E$. We denote by $E^\vee $ the graded dual $(E^\vee )_i = (E_{-i} )^*$, and $\mathcal{E} ^\vee $ its sheaf of sections.

There are a number of ways one could make precise the notion of functions on $\mathcal{E} (M)$. Of course, we are motivated to consider the symmetric algebra on the graded dual, but it will be helpful to consider other possibilities and to show their equivalence. One such possibility is the space of polynomials $P(\mathcal{E} (M))$. We intuitively expect the symmetric algebra of the dual to recover polynomials, so to see this intuitively, we start with a more intuitive definition of polynomials.

Costello and Gwilliam in \cite{cg1} start with a definition of polynomials in terms of differentiable vector spaces. I think our class didn't go over this machinery, so I will skip to their identification
\[
	P_n(\mathcal{E} (M)) \coloneqq \mathcal{D} _c\left(M^n , (E^!)^{\boxtimes n}\right)_{S_n}.
\]
Here $P_n$ denotes polynomials of degree $n$, $\mathcal{D} _c$ denotes compactly supported distributional sections of $(E^!)^{\boxtimes n}$, and $E\boxtimes E = E^{\boxtimes 2} $ denotes the second external tensor power of the bundle $E$, which is a $(\dim E)^2$-bundle over $M^2$ defined by, for $p,q\in M$,
\[
	E^{\boxtimes 2}_{(p,q)} =E_p\otimes E_q
.\]
We then take $S_n$-coinvariants. Here the compactly supported distributional sections of $M^n$ on $(E^!)^{\boxtimes n} $ are defined to be compactly supported continuous linear functionals on sections. For $n=1$, this looks like $f : \Gamma ^\infty (M,E) \to \mathbb{R} $. We do not go into detail about how this arises, and simply note that earlier we defined distributions as the dual to $C^\infty (M)$, which makes sense. Here we are trying to do something more general: define a dual to sections of an arbitrary vector bundle $E$. Apparently $\mathcal{D} (M,E)$ is defined as the linear functionals $\Gamma ^\infty _c(M,E^!) \to \mathbb{R} $, and so considering compactly supported sections $\mathcal{D}_c(M,E^!)$ dualizes this definition in a sense.

We expect $P(\mathcal{E} (M)) \coloneqq \bigoplus_{n} P_n(\mathcal{E} (M))$ to be the observables on a theory. However, there are a few other possible characterizations (we have already seen the derived critical locus). We continue going through them to show that they all make sense, and are all equivalent in this case. This will provide good intuition for how to formalize a free field theory axiomatically.

Since $E$ is graded, the evaluation map $\left\langle f,v \right\rangle=f(v) $ taking a distributional section $f$ and a section $v$ must respect grading in the sense that $\left\langle f,v \right\rangle $ has degree $\left| f \right| +\left| v \right| $. Since $\mathbb{R} $ is concentrated in degree 0, this means $f(v)$ is nonzero precisely when $\left| f \right| =-\left| v \right| $. This makes sense since the graded dual flips the direction. Hence $\mathcal{D} _c$ produces a whole complex in this case with $f : \Gamma (M,E_i) \to \mathbb{R} $ in degree $-i$.


We see fit to expand a bit on external tensor powers. We use a graded external tensor power
\[
	(E^!)^{\boxtimes n}_k = \bigoplus_{d_1+ \cdots + d_n=k} E^!_{d_1} \boxtimes \cdots \boxtimes E^!_{d_n} 
.\]
Compactly supported distributional sections of this bundle are compactly supported continuous linear functionals on sections $f : \Gamma ^\infty (M^n, E^{\boxtimes n} ) \to \mathbb{R} $. There is a canonical $S_n$-action by permuting the ordering of the base manifold. A section of $(E^!)^{\boxtimes n} $ can be represented as a tensor product of sections $s_1 \otimes \cdots \otimes s_n$ of $E^!$ in coordinates sufficiently close to a point. One simply permutes these by, for any $\sigma \in S_n$, mapping $s\mapsto \varepsilon (s,\sigma )s_{\sigma^{-1} (1)} \otimes \cdots \otimes s_{\sigma ^{-1} (n)} $ for $\varepsilon $ some Koszul sign rule. One then gets the dual action $\sigma \cdot f \coloneqq f\circ \sigma ^{-1} $, where the $\sigma ^{-1} $ acts on sections. Hence
\[
	(\sigma \cdot f)(s) = f(\sigma^{-1}  \cdot s)
.\]
Thus the $S_n$-coinvariants are symmetric linear functionals.

We explain why these are polynomials. Let $V$ a vector space. Then there is a basis $\left\{ e_i \right\} _{i\in I} $ of $V$, and a polynomial in $e_i$ can be thought of as a linear combination in formal powers of linear functions $\varepsilon^i(e_j)=\delta _{ij} $. The intuition is that, for example if we conside $f(x,y)=x^2+y$ and let $\chi (x)=1$ and $\eta (y)=1$ and otherwise zero, then to evaluate $(x,y)=(1,2)$, we evaluate
\[
	\chi ^2(1x+2y) + \eta (1x+2y) = \chi ^2(x) + 2\eta (y) = 1^2 + 2
.\]
Thus a polynomial can be thought of as an element of the symmetric algebra of the dual. This is just a graded generalization.

Now we return to our concrete example. Our derived space of fields can be realized as the sections of a graded vector bundle given by the trivial bundle $M\times \mathbb{R} $ concentrated in degrees 0 and 1. Write $E_0$ and $E_1$ for the 0th and 1st degrees of $E$, such that $E = E_0 \oplus E_1[-1]$ (we still make the degrees explicit). Then
\[
	E^! \coloneqq E^\vee \otimes \Lambda ^nT^*M = (E_0^* \oplus E_1^*[1])\otimes \Lambda ^nT^*M \cong (E_0^* \otimes \Lambda ^nT^*M) \oplus (E_1^* \otimes \Lambda ^nT^*M)[1]
.\]
The book claims that $E^!$ is just two copies of the trivial bundle, one in degree $-1$ and one in degree 0. Since $M$ is Riemannian, it has a canonical volume form, which yields a canonical basis for the 1-dimensional fibers of $\Lambda ^nT^*M$. Thus we can identify the fibers $(E_1^* \otimes \Lambda ^nT^*M)_p \cong (E_1^*)_p$ by the explicit isomorphism $f\otimes \mathrm{vol} _g \mapsto f$. Additionally, $\mathbb{R} $ has a canonical basis given by 1, so we can identify $\mathbb{R} ^*\cong \mathbb{R} $, inducing an isomorphism $(M\times \mathbb{R} )^* \cong M\times \mathbb{R} $ on the trivial bundles. This reconciles $E^!$ with the book's claims.

Note that since $P_1(\mathcal{E} (M))$ is just maps $\Gamma (M,E)\to \mathbb{R} $, and since as seen earlier it satisfies the same degree flipping as the graded dual, we obtain $P_1(\mathcal{E} (M))\cong \mathcal{E} (M)^\vee $:
\[
	\mathcal{E} (M)^\vee = \left( 
		\begin{tikzcd}
			\mathcal{D} _c(M)[1] \ar[" \Delta  ", r] & \mathcal{D} _c(M)
		\end{tikzcd}
	\right) 
.\]
The differential is still $\Delta $ since it is self-adjoint. Here we note that since monomials inherit a trivial $S_n$-action, passing to $S_n$-coinvariants is trivial, and the algebraic dual of $C^\infty (M)$ is just $\mathcal{D} _c(M)$.

One might think that an equivalent definition to a polynomial, and hence an observable, should be the symmetric algebra on the monomials $\mathcal{E} (M)^\vee $. This turns out to be a poor choice for quantization. However, free field theories offer a very nice classification of observables. Note that our identification of functions and densities yields an isomorphism
\[
	\mathcal{E}_c ^! (U)\cong \left( 
		\begin{tikzcd}
			C^\infty _c(U)[1]\ar[" \Delta  ", r] & C^\infty _c(U)
		\end{tikzcd}
	\right) 
\]
for compactly supported sections of $\mathcal{E} ^!$. Recall this is defined by taking a compactly supported section $f\otimes \mathrm{vol} _g$ of $(U\times \mathbb{R} )\oplus \Lambda ^nT^*U$ and identifying it with the compactly supported function $f : U \to \mathbb{R} $. This yields a natural map of cochain complexes  $\mathcal{E} ^!_c(M)\to \mathcal{E}^\vee $ by viewing a compactly supported function as a distribution (integrate against it). Here we note that the differentials on these complexes are the same; just trace definitions.

\begin{lemma}\label{lma:3.4.10}
	The map $\mathcal{E} _c^!(U)\to \mathcal{E} ^\vee (U)$ is a homotopy equivalence of cochain complexes.
\end{lemma}

The relevance of \cref{lma:3.4.10} is that we are now working \emph{homotopically}. As discussed earlier, we can recover the classical observables and the classical critical locus by taking cohomology of the relevant derived notions. Hence we expect observables to live in (the 0th) cohomology of some complex, meaning we only care about the derived notions up to homotopy. If observables in $\mathcal{E} (M)^\vee $ are closed distributional sections, then we can identify them with compactly supported smooth maps $f : U \to \mathbb{R} $ which are closed in $\mathcal{E} _c^!(U)$ in a homotopical sense. We think of the smooth observables as ``smeared''. For example, we intuitively replace the delta function $\delta_x$ at $x$ with a bump function supported near $x$.

This gives us a nice notion of observables: \emph{smeared observables}
\[
	\operatorname{Obs} ^{cl} (U)\coloneqq \operatorname{Sym} (\mathcal{E} _c^!(U)) = \operatorname{Sym} \left( 
		\begin{tikzcd}
			C_c^\infty (U)\ar[" \Delta  ", r] & C_c^\infty (U)
		\end{tikzcd}
	\right) 
\]
by taking the symmetric algebra, just as we expected to take the symmetric algebra on $\mathcal{E} ^\vee $. We claim that we may identify $\operatorname{Sym} ^n(\mathcal{E} _c^!(U))\cong C^\infty _c(U^n,(E^!)^{\otimes n} )_{S_n} $, which can be identified with a subspace of $P_n(\mathcal{E} (U))$ defined by taking all distributions to have compact support (why have we gone $\boxtimes \to \otimes $?). Here we take the graded symmetric algebra with respect to a topologically completed tensor product $\otimes $. This does not change the algebraic machinery, and simply completes the dense subset $C^\infty(M)\otimes C^\infty (N)\subseteq C^\infty (M\times N)$. To summarize,

\[
	\operatorname{Sym} (\mathcal{E} ^\vee (U))\simeq \operatorname{Obs} ^{cl} (U)\coloneqq \operatorname{Sym} (\mathcal{E} ^!_c(U)) \subseteq P_n(\mathcal{E} (U))
.\]


\begin{lemma}\label{lma:3.4.11}
	The map $\operatorname{Obs} ^{cl} (U)\to P(\mathcal{E} (U))$ is a homotopy equivalence of cochain complexes.
\end{lemma}
\begin{proof}
	It suffices to show the map $\operatorname{Sym} ^n(\mathcal{E} _c^!(U))\subseteq P_n(\mathcal{E} (U))$ is a homotopy equivalence for all $n$. For this, it suffices to show the inclusion $C^\infty _c(U,(E^!)^{\otimes n} ) \subseteq \mathcal{D} _c(U,(E^!)^{\otimes n} )$ is an $S_n$-equivariant homotopy equivalence. Finish later.
\end{proof}
Hence we have shown that all of our obvious notions of observables on this theory are equivalent, and distributional polynomial observables (given by integrating against some distribution on $U^n$) with smooth polynomial observables (given by integrating against a smooth function on $U^n$).

\end{example}


How do we intepret the observables we constructed in \cref{4.2.1}? We start by describing $\Obs^{cl} (U)$ more explicitly. Recall that $\mathcal{E} _!^c \cong \mathcal{E}_c^\vee = C^\infty_c \oplus C^\infty_c[-1]$. Thus thee complex looks like
\[
	\Sym (\mathcal{E} _c(U)) = \Sym (C^\infty _c(U) \oplus C^\infty_c(U)[-1])\cong \Sym C^\infty _c(U) \otimes \Lambda ^{\bullet} C^\infty _c(U),
\]
where $\Lambda ^{\bullet} C^\infty _c(U)$ is concentrated in negative degees, and $\Sym C^\infty _c(U)$ is concentrated in degree 0. Diagramatically,
\[\Obs^{cl} (U) \cong \left(\begin{tikzcd}[row sep = 2.5em, column sep = 2em]
\cdots \ar[r]  & \Lambda ^2 C^\infty _c(U)\otimes \operatorname{Sym} C^\infty _c(U)\ar[r]  & C^\infty _c(U)\otimes \Sym C^\infty _c(U)\ar[r]  & \Sym C^\infty _c(U)
\end{tikzcd}\right).\]
All tensor products here are topologically completed tensor products.

By identifying $C^\infty _c(U) \subseteq \mathcal{D}(U)=(C^\infty (U))^*$, we may think of $\Sym C^\infty _c(U)$ as an algebra of polymomial functions on $C^\infty(U)$, using the motivation we explored earlier. We can actually obtain similar intuition from the other terms in the complex. Let $T_cC^\infty (U) \subseteq TC^\infty (U)$ be the subbundle of the tangent bundle given by the subspace $C^\infty _c(U)$. ``An element of a fiber of $T_cC^\infty (U)$ is a first-order variation of a field which vanishes outside a compact set.'' Let's unpack that sentence from the authors of \cite{cg1}.

First, we view $T_cC^\infty (U)$ as a subbundle of $TC^\infty (U)$. Thus we are taking tangent vectors in the tangent bundle $TC^\infty (U)$, and restrictig them along $C^\infty _c(U)$. Since $C^\infty $ is a sheaf of vector spaces, $C^\infty (U)$ has trivial tangent bundle, so we obtain $TC^\infty (U)\cong C^\infty (U)\times C^\infty (U)$. Thus fibers of $TC^\infty (U)$ are pairs $(\phi ,\delta \phi )$ for $\phi \in C^\infty (U)$, and $\delta \phi $ a tangent vector to $\phi $, which we refer to as a \emph{first order variation} of $\phi $, and define via the calculus of variations.

Now to find elements of $T_cC^\infty (U)$, we restrict $\delta \phi $ along $C_c^\infty (U)$, meaning that $T_cC^\infty (U) \cong C^\infty (U)\times C_c^\infty (U)$. Thus elements of $T_cC^\infty (U)$ are compactly supported first order variations. Apparently this defines an integrable foliation on $T_cC^\infty (U)$, and this foliation can be defined for any sheaf of spaces.

We can interpret an element of $C_c^\infty (U)\otimes \Sym C^\infty _c(U)$ as a polynomial section of $T_cC^\infty (U)$: a section of $T_cC^\infty (U)$ is a map $X : C^\infty (U) \to C^\infty (U)\times C^\infty _c(U)$ which is a section of the projection; equivalently it is a map $X: C^\infty (U) \to C_c^\infty (U)\subseteq \mathcal{D} (U)$. A \emph{polyomial} in this sense means that component functions are polynomials, so a polynomial is an element of $C_c^\infty (U) \otimes \Sym \mathcal{D} (U)$ by defining $X(v) = \sum_{i} f_i(v)w_i$ for $f_i \in \Sym \mathcal{D} (U)$ and $w_i\in C_c^\infty (U)$. Since $C_c^\infty (U)\subseteq \mathcal{D} (U)$, we may interpret $C_c^\infty (U)\otimes \Sym C^\infty _c(U)$ as a space of polynomial sections of $T_cC^\infty (U)$. By the same logic, $\Lambda^kC_c^\infty (U)\otimes \Sym C^\infty _c(U)$ is the space of polynomial sections of $\Lambda ^kC_c^\infty (U)$.

In particular, if $\mathrm{PV}_c(C^\infty (U))$ denotes polynomial polyvector fields on $C^\infty (U)$ along the subbundle $T_cC^\infty (U)$, then
\[
	\Obs^{cl} (U)\cong \mathrm{PV} _c(C^\infty (U))
\]
as graded vector spaces.



\begin{remark}\label{4.2.4}
	We offer some motivating intuition for the tangent bundle to an infinite dimensional manifold. The definition of a tangent vector as a \emph{derivation} diverges from the notion of a geometric tangent vector in infinite dimensions, since there are strictly more derivations than geometric tangent vectors. Instead, we use the definition via equivalence classes of smooth curves. Recall that given $F : M \to N$ a smooth map between smooth manifolds, $p\in M$, and $v\in T_pM$, we find $\dd F_p(v) = (F\circ \gamma )'(0)$ for any open $0\in J\subseteq \mathbb{R} $ and any smooth $\gamma : J \to M$ with $\gamma (0)=p$ and $\gamma '(0)=v$. Hence one could \emph{define} a tangent vector at $p\in M$ as an equivalence class of smooth curves $\gamma : (-\varepsilon ,\varepsilon ) \to M$ with $\gamma (0)=0$, where we define $\gamma _1 \sim \gamma _2$ iff, for any chart $(p\in U, \varphi )$,
	\[
		\left. \frac{\mathrm{d}}{\mathrm{d}t} \right|_{t=0} (\varphi \circ \gamma _1)(t) = \left. \frac{\mathrm{d}}{\mathrm{d}t} \right|_{t=0} (\varphi \circ \gamma _2)(t)
	.\]
	The intuition is that this construction is just $\mathrm{d} \varphi_p(v)$, which is precisely the isomorphism $T_pM \cong \mathbb{R} ^{\dim M} $. Thus computing this with two different smooth curves should give the same result.
\end{remark}

We want to promote our identification $\Obs^{cl} (U) \cong \mathrm{PV} _c(C^\infty (U))$ to one of cochain complexes. Recall the action
\[
	S(\phi )=\frac{1}{2} \int_{U} \phi \Delta \phi 
\]
for $\Delta $ the Laplacian on the Riemannian manifold $(M,g)$. Roughly, the  differential on $\mathrm{PV}_c(C^\infty (U))$ should be given by contraction with $\mathrm{d} S$. We have to make this idea precise. Since $S$ is quadratic in $\phi $, we expect that contraction with a tangent vector $\phi _0\in C^\infty (U)$ should be the linear functional
\[
	(\phi _0\intprod \mathrm{d} S)(\phi )=\frac{1}{2} \int_{U} \phi _0 \Delta \phi 
.\]
However, recall that this is not well-defined for all fields $\phi $. However, it \emph{does} make sense to say
\[
	\frac{\partial S}{\partial \phi _0} (\phi )=\frac{1}{2} \int_{U} \phi _0 \Delta \phi 
\]
for all fields $\phi $, provided $\phi _0\in C^\infty _c(U)$ is compactly supported. More precisely, the desired form $\dd S$ doesn't make sense as a section of the tangent bundle, but it \emph{does} make sense as a section of $T_c^*C^\infty (U)$, the dual of $T_cC^\infty (U)$. We may thus contract with this 1-form on $\mathrm{PV} _c(C^\infty (U))$, showing that this differential indeed agrees with that for $\Obs^{cl} (U)$.

\textbf{Summary.} We identify $C_c\infty (U)\subseteq \mathcal{D} (U) = (C^\infty (U))^*$ by viewing compactly supported maps $f : U \to \mathbb{R} $ as functionals
\[
	g\mapsto \int_{U} (f\cdot g)\dd vol
.\]
We start with the sheaf of sections $\mathcal{E} $ of the derived space of solutions $E$, given by the sheaf of cochain complexes $\Delta : C^\infty (U)\to C^\infty (U)[-1]$. We expect the symmetric dual of this to be polynomial polyvector fields, which are our observables. Thus we should have $\Sym \mathcal{E} ^\vee \simeq \Obs^{cl} $ (we are working up to homotopy). We also expect (degree $n$) polynomials to be $\mathcal{D} _c\left( M^n, (E^!)^{\boxtimes n}  \right)_{S_n} $ -- that is, symmetric compactly supported distributional sections of the $n$-th exterior tensor power of $E^! = E^\vee \otimes \Omega _M^n$. Since $\Delta $ is self-adjoint and $S$ is quadratic in the field $\phi $, we find that $\mathcal{E} _c^! \cong \left(\Delta : C^\infty _c(U) \to C^\infty _c(U)\right)$ has
\[
	\mathcal{E} ^!_c(U) \simeq \mathcal{E} ^\vee (U)
\]
by simply viewing compactly supported sections as distributions. This yields
\[
	\Obs^{cl}(U) \cong \Sym \mathcal{E} ^!_c(U) \simeq \Sym E^\vee (U) \simeq P(\mathcal{E} (U))
.\]
Finally, we find that by viewing $\Obs^{cl} (U)$ as $\Sym \mathcal{E} ^!_c(U)$, we may identify it with compactly supported polyvector fields $\mathrm{PV} _c(C^\infty (U))$, and we may take the differential to be contraction with the action $\dd S$. Hence, we have shown that in the case of free field theories, provided we make the correct topological adjustments, the infinite-dimensional case echoes the finite dimensional case. 








\section{General free field theories}

\begin{definition}\label{dfn:4.3.1}
	Let $M$ be a manifold. A \emph{free field theory} on $M$ is the following data.
	\begin{enumerate}
		\item A graded vector bundle $E$ on $M$, whose sheaf of sections is denoted $\mathcal{E} $, and whose compactly supported sections are denoted $\mathcal{E} _c$.
		\item A morphism of sheaves of $\mathbb{R} $-vector spaces (meaning that it is not $C^\infty _M$-linear!) $\dd : \mathcal{E}  \to \mathcal{E} $ of cohomological degree 1, making $\mathcal{E} (U)$ into a sheaf of cochain complexes. This must also be a differential operator, meaning locally it looks like a differential operator.
		\item Let $E^! \coloneqq E^\vee \otimes \mathrm{Dens}_M$, which on Riemannian manifolds is $E^\vee \otimes \Lambda ^nT^*M$. Let $\mathcal{E} ^!$ be the sections of $E^!$, and let $\dd ^!$ be the differential on $\mathcal{E} ^!$ inherited from that of $\mathcal{E} $ as described in \cref{rmk:4.3.2}. Then there is a natural pairing between $\mathcal{E} _c(U)$ and $\mathcal{E} ^!(U)$, which is compatible with differentials.
		\item There is an isomorphism $E \cong E^![-1]$ which is compatible with differentials (I assume meaning the induced isomorphism of sheaves of graded vector spaces is an isomorphism of sheaves of cochain complexes), with the property that the induced pairing of cohomological degree $-1$ on each $\mathcal{E} _c(U)$ is graded anti-symmetric.
	\end{enumerate}
\end{definition}

The complex $\mathcal{E} (U)$ is the derived space of solutions to the equations of motion of the theory on an open subset $U$. The rest of the data ensures that we can make sense of this field theory in the same way as for the free scalar field theory in the previous section. Hence we should (I hope???) be able to identify $\Sym \mathcal{E} ^\vee \simeq \Sym \mathcal{E} _c^!(U)$ with the derived critical locus, and should be able to interpret it as polynomial polyvector fields on $\mathcal{E} _c^!(U)$ in some way.

\begin{terminology}\label{term:4.3.2}
	Let $M$ be a manifold. To specify a free field theory on $M$, we usualy simply writ ethe sheaf of sections $\mathcal{E} $ of the graded vector bundle $\pi : E \to M$, and will refer to the tuple $(M,\mathcal{E} )$ as a free field theory. This is an abuse of notation since it leaves the differential $\dd : \mathcal{E}  \to \mathcal{E} $ and the isomorphism $E\cong E^![-1]$ implicit, but is sufficient for our purposes.

	The manifold $M$ intuitively represents \emph{spacetime}. As such, we can take $M$ to be finite-dimensional, meaning the struggles of doing infinite dimensional differential geometry have been replaced with the language of sheaves of sections of vector bundles, and derived algebraic geometry. The sheaf of sections $\mathcal{E} $ is called the \emph{derived space of fields}. When paired with the differential $(\mathcal{E} ,\dd )$, we call it the \emph{derived space of solutions to the Euler-Lagrange equations}. In practice, we rarely distinguish between th e two notions, and use them interchangeably. The intuition is that we will take an action functional $S$, extract the Euler-Lagrange equations of motion, and then encode it into the complex $(\mathcal{E} ,\dd )$. For example, in the scalar field theory example, the action functional yielded the equations of motion $\Delta \phi =0$, which we encoded into $\mathcal{E} $ by having $H^0(\mathcal{E} ) = \Ker \Delta $ be the classical critical locus.
\end{terminology}



\begin{remark}\label{rmk:4.3.2}
	As noted previously, we assume that the graded vector bundle $E$ is finite rank, such that we avoid annoying congergence issues. The pairing $\left\langle -,- \right\rangle : \mathcal{E} _c\otimes_{C^\infty _M} \mathcal{E} ^!\to \mathbb{R} $ is defined as follows. Let $s\in \mathcal{E} _c(U)$, and $\alpha \otimes \omega \in \mathcal{E} _c^! = \mathcal{E} (U)^\vee \otimes_{C^\infty_M(U)} \Omega _M^n(U)$, where we assume $M=(M,g)$ is Riemmanian to make the presentation more simple. There are no changes in the fully general case. Then we have that $\alpha =\left\{ \alpha ^i : \mathcal{E} (U)^{-i}  \to \mathbb{R}  \right\}_{i\in\mathbb{Z} } $.

	The direct sum perspective of the total space of a graded vector bundle is very useful here. Recall that $\mathcal{E} ^\vee =\Gamma (-,E^\vee )$, and $E^\vee $ is defined fiberwise by
	\[
		E^\vee_x=\HOM_{\mathbb{R} }(E^{-i}_x, \mathbb{R} ) 
	.\]
	The direct sum perspective implies that
	\[
		\alpha_x \in \bigoplus_{i\in\mathbb{Z} } (E^{-i}_x)^* 
	.\]
	By writing this as an element of a direct sum, we conclude that $\alpha $ can be written fiberwise as a linear combination of maps
	\[
		\alpha_x =\sum_{i\in\mathbb{Z} } \left( \alpha ^i_x : E^{-i}_x \to \mathbb{R}  \right) 
	,\]
	where of course $\alpha ^i_x\in (E^{-i} _x)^*$. Thus we can act an element $\alpha $ of the graded dual $\mathcal{E} ^\vee (U)$ on a section $s\in \mathcal{E} (U)$ to obtain a pairing $\mathrm{ev} : \mathcal{E} ^\vee \otimes_{\mathbb{R} } \mathcal{E}  \to \mathbb{R} $:
	\[
		\mathrm{ev} (\alpha ,s) = \alpha (s) = \sum_{i\in\mathbb{Z} } \alpha ^{-i} (s^i)
	.\]
	Now we can define the pairing. If $s$ is compactly supported, then $\alpha (s)$ is a compatly supported map $M\to \mathbb{R} $. This is easy to show by simply using $\alpha ^i_x(0)=0$. Thus, $\alpha (s)$ defines a derivation $(\alpha (s))^?$ by integrating against it. We have to integrate $n$-forms, so this determines a map $\Omega _M^n(U)\to \mathbb{R} $. This map is $C^\infty _M(U)$-linear:
	\[
		(\alpha (s))^?(f\omega) = \int_{U} \alpha (s)f\omega =(f\cdot \alpha (s))^?(\omega )
	.\]
	Transposing under $-\otimes_{C^\infty _M}V \dashv \Hom_{C^\infty _M}(V, -) $ yields the pairing
	\[
		\left\langle s,\alpha \otimes \omega  \right\rangle_U = \int_{U} \alpha (s)\omega =\int_{U} \sum_{i\in\mathbb{Z} } \alpha ^{-i} (s^i)\omega 
	.\]
	
	Now we show that this pairing is indeed a valid morphism in the category $\Fun(\Open(M) , \mathcal{G} \Mod_{C^\infty _M} ) $. That is, we must show that for all open $U\subseteq M$, the map $\left\langle -,- \right\rangle _U$ can be defined degreewise as a map $\mathcal{E} _c\otimes _{C^\infty _M(U)} \mathcal{E} ^! \to \mathbb{R} [0]$. Indeed, recall that if $\alpha ^{-j} \in \mathcal{E} ^\vee (U)^{-j} $ and $s^i\in \mathcal{E} (U)^i$, then by definition of the graded dual, $\alpha ^{-j} (s^i) = \delta _{ij} \alpha ^{-i} (s^i)$, for $\delta _{ij} $ the Kronecker delta. Thus the map vanishes unless $j=i$, meaning we obtain a map $\mathcal{E} _c^i\otimes _{C^\infty _M(U)} (\mathcal{E} ^!)^{-i} \to \mathbb{R}$. Thus this map respects degrees.


	
	We now define $\dd ^!$ as the (graded) adjoint of $\dd $ with respect to this pairing. That is, $\dd ^!$ is the unique map $\mathcal{E} ^! \to \mathcal{E} ^!$ with the property that
	\[
		\left\langle \dd s,\sigma  \right\rangle +(-1)^{\left| s \right| } (s,\dd^! \sigma )=0
	.\]
	We can obtain an explicit description of $\dd ^!$. Start by assuming that $s$ has even degree, so we can ignore the factor $(-1)^{\left| s \right| } $. Let $U\subseteq M$ be open, and let $s\in \mathcal{E} (U)^{i-1} $ and $\alpha\in \mathcal{E} ^{\vee } (U)$. Then since $\dd $ has cohomological degree 1, $\alpha (\dd s)$ is defined. Thus we may consider the map
	\[
		\left\langle \dd s,\alpha \otimes \omega  \right\rangle =\int_{U} \alpha (\dd s)\omega 
	\]
	for a volume form $\omega $. To explicitly describe $\dd ^!$, note that the pairing can be described as
	\[
		\left\langle -,-  \right\rangle =\int_{U} m\circ (\mathrm{ev} \otimes 1)
	,\]
	where $m : C^\infty _M(U)\otimes_{C^\infty_M(U)} \Omega _M^n(U) \to \Omega ^n_M(U)$ is the multiplication map, $\mathrm{ev} : \mathcal{E}^\vee (U)^{-i} \otimes_\mathbb{R} \mathcal{E} _c^i \to C^\infty _M(U)$ is the evaluation map, and 1 is the identity map.

	If we take $\dd ^!$ to be such that $\dd ^!_U(\alpha \otimes \omega )$ is the map $s\mapsto (\alpha (\dd s)\otimes \omega )$, then $\dd $ and $\dd ^!$ are necessarily adjoints when passing to the integral. We want to simply define $\dd ^!\coloneqq \dd ^*\otimes 1$, for $(-)^*$ precomopsition. If we do this, then
	\[
		\dd ^!(\alpha \otimes \omega ) = (\alpha \circ \dd )\otimes \omega 
	.\]
	The pairing is thus (up to a factor of $\pm$ depending on the degree):
	\[
		\left\langle s,\dd^!(\alpha \otimes \omega ) \right\rangle =\int_{U} (\alpha \circ \dd )(s)\omega = \alpha (\dd s)\omega =\left\langle \dd s,\alpha \otimes \omega  \right\rangle 
	.\]
	Unfortunately, $\alpha \circ \dd $ need not be a section of the graded dual vector bundle, so this is in general ill-defined. However, a formal adjoint always exists, meaning we can always find an operator which equates to $\dd ^* \otimes 1$ when passing to the integral. Thus we have the best description of our map that we can obtain without passing to examples. Moreover, note that $\dd ^!(\alpha \otimes \omega ): \mathcal{E} ^{i-1}  \to \mathbb{R} $, and thus this is an element of homogeneous degree $-(i-1) = -i+1$. Thus $\dd ^! : \mathcal{E} ^! \to \mathcal{E} ^![1]$.


	Consider the scalar field theory. The graded vector bundle $E$ was the trivial bundle ocncentrated in degrees 0 and 1, such that $\mathcal{E} $ was simply $C^\infty_M$ concentrated in degrees 0 and 1. The differential $\dd : C^\infty _M \to C^\infty _M$ was $\Delta +m^2$, for $\Delta $ the Laplacian on the Riemannian manifold $(M,g)$.

	Note that we can identify $(\mathbb{R} \times M)^* = \mathbb{R} ^*\otimes M$ with $\mathbb{R} \otimes M$ using the canonical basis for $\mathbb{R} $ given by the element 1. Thus we obtain
	\[
		\mathcal{E} ^\vee \cong \left( 
			\begin{tikzcd}
				C^\infty _M[1]\ar[" \Delta +m^2 ", r] & C^\infty _M
			\end{tikzcd}
		\right) 
	.\]
	The isomorphism $\mathbb{R} ^*\times M \cong \mathbb{R} \times M$ is given by evaluation, meaning a section $\alpha $ of $\mathbb{R} ^*\times M$ is identified with a section of $\mathbb{R} \times M$ by evaluating fiberwise:
	\[
		\alpha _x \mapsto \alpha _x(1)
	.\]
	Now note that $\Omega ^n_M\cong \operatorname{span}_{C^\infty _M} (\mathrm{vol} _g)$. Thus we have another canonical identification $C^\infty _M \otimes_{C^\infty _M} \Omega ^n_M$ by simply mapping $f\mapsto f\otimes \mathrm{vol} _g$. This yields
	\[
		\mathcal{E} ^! \cong \left( 
			\begin{tikzcd}
				C^\infty _M[1]\ar[" \Delta +m^2 ", r] & C^\infty (M)
			\end{tikzcd}
		\right) 
	\]
	where we identify $f\in C^\infty _M(U)$ with $\alpha \otimes \omega \in \mathcal{D} (U)\otimes \Omega _M^n(U)$ by defining $\omega =\mathrm{vol} _g$, and since $\alpha _x(1) = f(x)$, we identify $\alpha (g) = f\cdot g$. Thus the pairing is
	\[
		\left\langle f,g \right\rangle_U =\int_{U} fg\mathrm{vol} _g
	.\]
	We then recover that $\dd ^!$ is also $\Delta $ since the Laplacian is self-adjoint.
\end{remark}

\begin{remark}\label{rmk:4.3.4}
	Condition (4) of \cref{dfn:4.3.1} yields an isomorphism $E\cong E^![-1]$, which necessarily induces an isomorphism $\varphi : \mathcal{E} \cong \mathcal{E} ^![-1]$ on sheaves of sections. We can define a canonical pairing $\omega : \mathcal{E} _c\otimes_{C^\infty _M} \mathcal{E}_c \to \mathbb{R} [0]$ is given by
	\[
		\omega (s^i,s^j) \coloneqq \left\langle s^i,\varphi (s^j) \right\rangle 
	.\]
	We investigate what this means. First, we assume that $s^i\in \mathcal{E} (U)^i$ and $s^j\in \mathcal{E} (U)^j$. Then $\varphi (s^j)\in (\mathcal{E} ^!(U))^{j-1} $. Thus we can write $\varphi (s^j) = \alpha ^{-j+1} \otimes \omega $, where $\alpha ^{-j+1} : \mathcal{E} (U)^{-j+1}  \to \mathbb{R} $. Hence we see that $\omega (s^i,s^j)$ is going to vanish unless $-j+1=i$. Hence we can write
	\[
		\omega (s^{i} ,s^{-i+1} ) = \left\langle s^i,\varphi (s^{-i+1} ) \right\rangle 
	.\]
	We also denote this by $\left\langle \left\langle -,- \right\rangle  \right\rangle $ sometimes. We note that this is a map of cohomological degree $-1$, since $s^i \otimes s^{-i+1} $ has degree 1, and it maps into $\mathbb{R} [0]$.

	Condition (4) also dictates that
	\[
		\omega (s^i,s^j) = (-1)^{\left| s^i \right| \left| s^j \right| } \omega (s^j,s^i)
	.\]
	Applying this to the case where $\omega $ is nonzero:
	\[
		\omega (s^i,s^{-i+1} ) = -(-1)^{i(-i+1)} \omega (s^{-i+1} ,s^i) = -(-1)^{-i^2 + i} \omega (s^{-i+1} ,s^i) 
	\]
	\[
		=-(-1)^{-i+i} \omega (s^{-i+1} ,s^i) = \omega (s^{-i-1} ,s^i)
	.\]
	Thus ``graded antisymmetric'' in this case means ''symmetric in the normal way that you would normally think of it''.
\end{remark}

\begin{remark}\label{4.3.5}
	We remark on the actual meaning of the conditions in \cref{dfn:4.3.1}. Conditions (1) and (2) simply set up a system of linear differential equations, albeit using the advanced machinery of sheaf theory and homological algebra. However, they do not inherently carry any physical meaning. Sure, defining the differential to encode the equations of motion can write a physical theory as such a mathematical object, but the converse does not follow. We need more conditions in order for observables to make sense. This is what condition (4) does: it provides an identification that is necessary to obtain physical meaning. It is difficult to intuitively explain why this identification exists at this time, however. The thing that makes the theory \emph{free} is the fact that the sheaf $\mathcal{E} $ is valued in cochain complexes of \emph{vector spaces}. Interacting theories are usually modeled in $L_\infty $-algebras.
\end{remark}


For the next example, recall the following.

\begin{definition}\label{dfn:4.3.6}
	Let $\pi : E \to M$ be a vector bundle over a manifold $M$. A \emph{connection} on $E$ is an $\mathbb{R} $-linear map $\nabla : \Gamma (M,E) \to \Gamma (M,T^*M\otimes_\mathbb{R} E)$ which satisfies a product rule
	\[
		\nabla (fs) = \dd f\otimes s + f\nabla s
	\]
	for all $f\in C^\infty (M)$ and $s\in\Gamma (M,E)$.
\end{definition}

\begin{definition}\label{dfn:4.3.7}
	Let $(M,g)$ be a pseudo-Riemmanian $n$-manifold. Then for all $k$, the \emph{Hodge dual} of a $k$-form $\zeta $ is the $(n-k)$-form $\star \zeta $ which satisfies the property
	\[
		\eta \wedge \star \zeta =g(\eta ,\zeta )\mathrm{vol} _g
	\]
	for all $k$-forms $\eta$.
\end{definition}
Here we define $g(\eta ,\zeta )$ in local coordinates as $\sum_{i,j} g^{ij} \eta _i\zeta _j$, where $g^{ij} $ are the matrix components of the inverse matrix of $g_{ij} =g(\partial_i,\partial _j)$.

Apparently we can use the Hodge star to define the \emph{coderivative} $\Omega ^k_M \to \Omega ^{k-1} _{M} $, given by $\delta =(-1)^{k} \star ^{-1} \dd \star $. The relevance of this map is that if $\left\langle \left\langle \eta ,\zeta  \right\rangle  \right\rangle  $ is the $L^2$ inner product
\[
	\left\langle \left\langle \eta ,\zeta  \right\rangle  \right\rangle =\int_{U} \eta \wedge \star \zeta 
;\]
here $\eta ,\zeta \in \Omega ^k_M(U)$; then $\delta $ is the right adjoint of $\dd $ under this inner product. Then the \emph{Hodge Laplacian} $\dd \star \dd $ is defined as $\dd \star \dd \coloneqq \dd \delta +\delta \dd  = \left( \dd +\delta  \right) ^2$.



\begin{example}\label{ex:4.3.8}
	We give an example of how to encode gauge symmetry into a free field theory. We describe Abelian Yang-Mills theory with gauge group $\mathbb{R} $. Let $(M,g)$ a Riemannian 4-manifold. The ``naive fields'' are 1-forms on $M$ ($A\in \Omega ^1_M$), which we think of as connections on the trivial line bundle $\mathbb{R} \times  M$ via using $\mathrm{d} +A$ as the connection corresponding to the 1-form $A$.

	This is indeed a connection:
	\[
		(\dd +A)(f) = \dd f + Af;
	\]
	here we use the canonical isomorphism $TM\cong T^*M$ for Riemannian manifolds to define $Af$ for 0-forms $f$. We then have $T^*M \otimes _\mathbb{R} (M\times \mathbb{R} ) \cong T^*M$, so the product rule is immediate or something idk.

	The abelian Yang-Mills action is
	\[
		S_{\mathrm{YM} } (A)=-\frac{1}{2} \int_{M} \mathrm{d} A\wedge \star \mathrm{d} A=\frac{1}{2} \int_{M} A(\mathrm{d} \star \mathrm{d} )A,
	\]
	where $\star \mathrm{d} A$ is the Hodge star on $\mathrm{d} A$, and $\mathrm{d} \star \mathrm{d} $ is the \emph{Hodge Laplacian on 1-forms (up to another factor of $\star $)}. We can prove this equality:
	\[
		\dd (A \wedge \star \dd A) = \dd A \wedge \star \dd A - A \wedge \dd (\star \dd A)
	.\]
	Integrating both sides, and applying Stokes' theorem to the left, and also assuming that $M$ has no boundary yields
	\[
		\int_{M} \dd A \wedge \star \dd A = \int_{M} A \wedge \dd (\star \dd )A
	.\]
	The notation $A(\dd \star \dd )A$ is an abuse of notation for $A\wedge \dd (\star \dd A)$.

	The Euler-Lagrange equations here are that $(\mathrm{d} \star \mathrm{d})A=0$. The derivation here is mildly annoying, and we omit it. This theory has a symmetry, which is that if $X\in\Omega ^0_M$, then $A\mapsto A+\mathrm{d} X$. This is a symmetry because $\mathrm{d} ^2X=0$, so plugging in $A+\mathrm{d} X$ into the action yields
	\[
		S_{\mathrm{YM} } (A+\mathrm{d} X)=S_{\mathrm{YM} } (A)
	.\]
	Hence these configurations would be considered \emph{physically equivalent}, yielding a group of symmetries on the na\"ive fields of our theory.

	The derived space of fields representing this situation is given by the complex
	\[\begin{tikzcd}[row sep = 2.5em, column sep = 2.5em]
		\mathcal{E} : \Omega ^0(M)[1]\ar[" \mathrm{d}  ", r]  & \Omega ^1(M)\ar[" \mathrm{d} \star \mathrm{d}  ", r]  & \Omega ^3(M)[-1]\ar[" \mathrm{d}  ", r] & \Omega ^4(M)[-2]
	\end{tikzcd}\]
	We explain why. Note that similarly to the scalar field theory example, the map $\dd \star \dd $ going from degree 0 to degree 1 cuts out the classical critical locus as its kernel $\Ker (\dd \star \dd )$. Hence if there was no degree $-1$ component (as there was in the scalar field theory), we could recover the classical space of solutions as $H^0(\mathcal{E} )$. This justifies the term ``derived space of solutions'' or ``derived space of fields'' (which we use more or less interchangeably).

	However, we also have our collection of symmetries which yield physically equivalent theories. Hence we should set $A$ equal to $A+\dd X$ when passing to cohomology. The way we do this is by adding a degree $-1$ component with image in $\Omega ^1(M)$ given by $\dd \Omega ^0(M)$, which is exactly what we see in the derived space of fields. Thus we obtain
	\[
		H^0(\mathcal{E} )=\frac{\text{Fields} }{\text{Gauge symmetries} } 
	.\]
	The last component will not be explained yet.
\end{example}

The idea we obtain from \cref{ex:4.3.8} is that classical fields can be recovered as the homology of the derived space of fields. Additionally, the derived space of fields gives us the necessary machinery required to encode the notion of a \emph{Gauge symmetry} into our field theory.

\begin{definition}\label{dfn:4.3.9}
	Let $M$ have a free field theory on it with derived space of fields $\mathcal{E} $. Then the \emph{classical observables} of the theory are defined as
	\[
		\Obs^{cl} (U) = \Sym \left( \mathcal{E} _c^!(U) \right) \cong \Sym \left( \mathcal{E} _c(U)[1] \right) 
	.\]
\end{definition}

We explore \cref{dfn:4.3.9}. First, we explore the isomorphism. If $(M,g)$ is Riemannian, then
	\[
		\mathcal{E} ^!_c = \Gamma ^\infty _c(-,E^\vee \otimes_\mathbb{R} \Lambda ^nT^*M) \cong \Gamma ^\infty _c(-,E^\vee )\otimes_{C^\infty _M} \Gamma ^\infty _c(-,\Lambda ^nT^*M) = \mathcal{E} ^\vee _c \otimes _{C^\infty _M} (\Omega ^n_M)_c
	.\]
	
Our intuition follows from the scalar free field theory example. The main reason that we were able to identify $\mathcal{E} _c^!$ with $\mathcal{E} ^\vee (U)$ (which we expect to be observables) up to quasi-isomorphism is the isomorphism $\mathcal{E} ^! \cong \mathcal{E} [-1]$. Since this aspect of a free field theory is part of the definition, we always have a notion of smeared observables on a free field theory.


\begin{lemma}\label{lma:4.3.10}
	Let $(M,\mathcal{E} )$ be a free field theory on a manifold $M$. Then the classical observables $\Obs^{cl} $ form a prefactorization algebra valued in commutative differential graded $\mathbb{R} $-algebras.
\end{lemma}
\begin{proof}
	First, we recall that the symmetric algebra is by definition a commutative graded algebra. Additionally, the differential $\dd ^! : \mathcal{E} ^! \to \mathcal{E} ^!$ extends immediately by simply defining
	\[
		\dd^!(ab) = (\dd^!a)b + (-1)^{\left| a \right| } a(\dd b)
	\]
	for all homogeneous elements $a,b$. Thus $\Sym \mathcal{E} ^!_c$ is a precosheaf of commutative differential graded algebras. To be explicit, we describe the algebra homomorphisms $i^U_V : \Obs^{cl} (U) \to \Obs^{cl} (V)$ induced by inclusions $U\subseteq V$ of open sets. Let $s\in \mathcal{E} _c^!(U)$. Then one can define $s|(V\setminus U) \coloneqq 0$, and the extension is still compactly supported. One then extends this to the entire symmetric algebra by simply defining
	\[
		i^U_V(st) = i^U_V(s)i^U_V(t)
	.\]
	This is immediately seen to respect the degrees, so we need only show that it respects the differentials.

By the universal property of the symmetric algebra, it suffices to show this on the generators $\mathcal{E} _c^!(U)$. Let $s\in \mathcal{E} _c^!(U)^i$; then $\dd ^! s\in \mathcal{E} _c^!(U)^{i+1} $. Since extension morphisms act trivially on $s$ on the original domain (in the sense that $(i^U_Vs)|U=s$), it suffices to show that $\dd ^!i^U_Vs$ and $i^U_V\dd ^!s$ agree on $V\setminus U$. Since $(i^U_V\dd ^!s)|(V\setminus U)=0$, it suffices to show that $(\dd ^!i^U_Vs)|(V\setminus U)=0$.

We use the fact that $\dd ^!$ is a morphism of sheaves, meaning it commutes with the restriction maps. In particular, suppose that there is some open set $W\subseteq V$ such that $s|W=0$. Then since the empty set is compact, $s|W$ is compactly supported, and moreovoer $(s|W)_x=s_x$ by the explicit structure of $\mathcal{E} ^!$. We then obtain for any $x\in W$,
	\[
		(\dd ^!s)_x = ((\dd ^!s)|W)_x = (\dd ^!(s|W))_x = (\dd ^!0)_x=0
	.\]
	Thus if $s|W=0$, then $(d^!s)|W=0$.

	Although $V\setminus U$ need not be an open set, we know that $V\setminus \operatorname{supp} s = V \cap  (\operatorname{supp} s) ^{c} $ is a finite intersection of open sets, and is thus open. Thus we can take $W = V\setminus \operatorname{supp} s = V\setminus \operatorname{supp} (i^U_Vs)$, and use the fact that $(i^U_Vs)|W=0$ to obtain that $(\dd ^!i^U_Vs)|W=0$. Since $\operatorname{supp} s \subseteq U$, we obtain $V\setminus U \subseteq W$, proving the statement. Therefore, $\Obs^{cl} $ is a precosheaf valued in commutative differential graded algebras.

	Now we prove $\Obs^{cl} $ is a prefactorization algebra. We recall that this means for every finite collection of pairwise disjoint open sets $\left\{ U_i \right\}_{1\leq i\leq n} $ and every open set $V$ containing $U_i$ for all $i$, there is a structure map $\bigotimes_{i} \Obs^{cl} (U_i) \to \Obs^{cl} (V)$ which satisfies the obvious coherence axiom. We define this map by defining the $n$-multilinear map
	\[\begin{tikzcd}[row sep = 0em, column sep = 2.5em]
		\displaystyle{\prod_{1\leq i\leq n} \Obs^{cl} (U_i)}\ar[r]  & \Obs^{cl} (V) \\
		(\alpha _1, \ldots ,\alpha _n)\ar[r, mapsto]  & \displaystyle{\prod_{1\leq i\leq n} i_{V}^{U_i} \alpha _i,}
	\end{tikzcd}\]
	where the product $\prod_{i} i_V^{U_i}\alpha _i$ denotes the symmetric product in the symmetric algebra. Here we take each $\alpha _i$ to be a homogeneous element of $\Obs^{cl} (U)$ of arbitrary (but not necessarily the same!) degree. Universality allows this to descend to a map on the graded tensor product. Additionally, by definition of the differential on the product of cochain complexes, and by definition of the differential on the symmetric algebra, this is a morphism of cochain complexes, meaning the induced map is also a morphism of cochain complexes.

	It is important to note that on fields $\phi \in \mathcal{E}(V)$, we define the action of $\prod_{i} i_V^{U_i}\alpha _i$ on $\phi $ as follows:
	\[
		\left( \prod_{1\leq i\leq n} i_{V} ^{U_i} \alpha _i \right) (\phi )=\prod_{1\leq i\leq n} \alpha _i(\phi |U_i),
	\]
	where the product on the right is usual multiplication in $\mathbb{R} $. This makes physical sense, since one expects the product of the observables to be the... product of the observables. The coherence axiom is immediately checked by simply writing out the definitions. Hence classical observables form a prefactorization algebra valued in cochain complexes.
\end{proof}
	





\section{The one-dimensional case, in detail}


We illustrate an example computing $H^0(\operatorname{Obs} ^{cl} )$ for a 1-dimensional free field theory on $\mathbb{R} $. We note that physically, this theory should reflect mechanics of a single particle in 1 dimension. Let $U\subseteq \mathbb{R} $ be an open subset. The action functional is
\[
	S_U(\phi )=\int_{U} \phi (\Delta+m^2) \phi \dd x
.\]
Here we note that $\mathrm{vol} _g=\dd x$. We can write this in another form using interation by parts:
\[
	\int_{U} \phi \partial _x^2\phi \dd x = \left[\phi \partial _x\phi\right]_{\partial U}  - \int_{U} (\partial _x\phi)(\partial _x\phi )\dd x
.\]
Recall that the action functional is generally ill-defined, and to recover mathematical intuition from it we use compactly supported sctions to extract the Euler-Lagrange equations. Hence for all mathematical purposes, we can take this to vanish on $\partial U$, meaning we can also take
\[
	S_U(\phi ) = \int_{U}\left( (\partial _x\phi )^2 + m^2\phi ^2\right)\dd x
.\]

Our first important result in this section is the following.
\begin{lemma}\label{lma:4.4.1}
	If $U=(a,b)\subseteq \mathbb{R} $ is an open interval, then
	\[
		H^{\bullet} (\Obs^{cl} (U))\cong \mathbb{R} [p,q]
	\]
	is the polynomial algebra on two generators.
\end{lemma}
As we can see, we can interpret $p$ as position and $q$ as momentum, which makes physical sence since all observables can be written using position and momentum.
\begin{proof}
Recall that the derived space of fields for the scalar free field theory
\[\mathcal{E}(a,b) = \left(\begin{tikzcd}[row sep = 2.5em, column sep = 2.5em]
	C^\infty (a,b)\ar[" \Delta +m^2 ", r]  & C^\infty (a,b)[-1]
\end{tikzcd}\right).\]
Recall that up to quasi-isomorphism, we can identify
\[\mathcal{E} _c^!(a,b) = \left(\begin{tikzcd}[row sep = 2.5em, column sep = 2.5em]
	C^\infty_c(a,b)[1]\ar[" \Delta +m^2 ", r]  & C^\infty _c(a,b)
\end{tikzcd}\right).\]
We will show that the second complex is quasi-isomorphic to $\mathbb{R} ^2$ concentrated in degree 0. The result then follows.

In the book, we use the convention
\[
\Delta =-\left( \frac{\partial }{\partial x}  \right) ^2
.\]
Thus the equations of motion say
\[
	\left( -\frac{\partial ^2}{\partial x^2} +m^2 \right) \phi =0
.\]
We can solve this differential equation because it is very easy. The solutions depend on if $m=0$, $m>0$, or $m<0$. If $m=0$, then $\phi $ is a linear function, so the solutions are spanned by $1$ and $x$. If $m>0$, then we want
\[
	\frac{\partial ^2\phi }{\partial x^2} =m^2\phi 
.\]
The solutions are spanned by $e^{\pm mx} $. Similarly, if $m<0$, then they are spanned by $e^{\pm i mx} $.

We use a nice basis here. If $m=0$, then define $\phi _q(x)=1$ and $\phi _p(x)=x$. If $m>0$, then choose the basis
\[
	\phi _p(x)=\frac{1}{2m} \left( e^{mx} -e^{-mx}  \right) ,\qquad \phi _q(x)=\frac{1}{2} \left( e^{mx} +e^{-mx}  \right) 
.\]
The point of this choice is that the adjoint of the relations between position and momentum that are satisfied in physics are precisely satisied by $\phi _p$ as position and $\phi _q$ as momentum. For example,
\[
	\phi _p(0)=1,\qquad \phi_q(0)=0,\qquad \phi _q'(0)=1,\qquad \phi _p'=\phi _q.
\]
Without loss of generality, assume $(a,b)=(-1,1)$. Define $\pi : \mathcal{E} _c(-1,1) \to \operatorname{Span} (p,q)\eqqcolon \mathbb{R} \left\{ p,q \right\} $ by
\[
	g\mapsto \pi (g)=\left(\int_{\mathbb{R} } g\phi _q\mathrm{d} x\right)q + \left( \int_{\mathbb{R} } g\phi _p\mathrm{d} x \right) p
.\]
We claim that this defines a cochain map
\[\begin{tikzcd}[row sep = 2.5em, column sep = 2.5em]
	C_c^\infty (-1,1)\ar[" \Delta +m^2 ", r] \ar[" \pi ^{-1}  "', d]  & C^\infty _c(-1,1)\ar[" \pi ^0 =\pi ", d]  \\
	0\ar[r]  & \mathbb{R} ^2\left\{ p,q \right\} .
\end{tikzcd}\]
Indeed, we can use the fact that $\Delta +m^2$ is self adjoint to show that
\[
	\pi((\Delta +m^2)(g))=q\int_{\mathbb{R} } (\Delta +m^2)g\phi _q\mathrm{d} x + p \int_{\mathbb{R} } (\Delta +m^2)g\phi _p\mathrm{d} x 
\]
\[
	= q \int_{\mathbb{R} } g(\Delta +m^2)\phi _q\mathrm{d} x + p \int_{\mathbb{R} } g(\Delta +m^2)\phi _p\mathrm{d} x
.\]
Since $\phi _p$ and $\phi _q$ solve $(\Delta +m^2)\phi =0$, this is necessarily zero. Therefore this is a map of cochains. We claim that $\pi $ is a quasi-isomorphism if $\Ker \pi $ is acyclic, meaning $H^{\bullet} (\Ker \pi )\cong 0$. To see this, write the cochain complex $0\to \mathbb{R} ^2\left\{ p,q \right\} $ as $D$, and then surjectivity of $\pi $ necessarily yields a short exact sequence of cochain complexes
\[\begin{tikzcd}[row sep = 2.5em, column sep = 2.5em]
	0\ar[r]  & \Ker \pi \ar[r]  & \mathcal{E} _c\ar[r]  & D\ar[r]  & 0.
\end{tikzcd}\]
Taking cohomology yields the long exact sequence
\[\begin{tikzcd}[row sep = 2.5em, column sep = 1.8em]
	0\ar[r]  & H^{-1} (\Ker \pi ) \ar[r] & H^{-1} (\mathcal{E} _c)\ar[r]  & H^{-1} (D)\ar[r]  & H^0(\Ker \pi )\ar[r]  & H^0(\mathcal{E} _c)\ar[r]  & H^0(D)\ar[r]  & 0.
\end{tikzcd}\]
If we assume that $H^{\bullet} (\Ker \pi )\cong 0$, then this reduces to the two exact sequences
\[\begin{tikzcd}[row sep = 0em, column sep = 2.5em]
	0\ar[r]  & H^{-1} (\mathcal{E} _c)\ar[r]  & H^{-1} (D)\ar[r]  & 0 \\
	0\ar[r]  & H^0(\mathcal{E} _c)\ar[r]  & H^0(D)\ar[r]  & 0,
\end{tikzcd}\]
and exactness implies these are isomorphisms. The definition of the long exact sequence tells us that these maps are precisely the homology of $\pi $, meaning that $\pi $ lowers to an isomorphism in cohomology, and is thus a quasi-isomorphism. It suffices to show this.

It suffices to show that $\Ker \pi $ is contractible. Recall that if $\alpha ,\beta : A \to B$ are cochain maps, then a homotopy $h: \alpha  \Rightarrow \beta $ is a cochain map $h : A \to B[-1]$ such that $f-g=d_B\circ h + h\circ d_A$. In our case, we require that $1-0=d_\mathcal{E} \circ h + h \circ d_{\Ker \pi } $ for 1 the identity. The map $\Ker \pi ^{-1} \to \Ker \pi ^0$ is just the restriction of $\Delta+m^2$. Note that since $h$ goes backwards in degree, the only $h$ that can be nonzero is the map $\Ker \pi ^0 \to C^\infty _c(-1,1)[1]$. Thus we simply need a map $h : C^\infty _c(-1,1) \to C^\infty (-1,1)$ which takes degree 0 elements to degree $-1$ elements, and it must satisfy $1 = d_{\mathcal{E} } \circ h$, meaning
\[
	(\Delta +m^2)h(g)=g
.\]

We recall the theory of Green's functions. Let $L_x$ be a differential operator on a manifold $M$ and $f : M \to \mathbb{R} $ compactly supported. A \emph{Green's function} is a distribution (not a function because it could have singularities, so it only makes sense to integrate against compactly supported forms) $G : M\times M \to M$ with $L_xG(x,y)=\delta (x-y)$. What we mean by this, at least in the case of $M=\mathbb{R} $, is as follows.
\[
	L\int_{\mathbb{R} } G(x,y)g\mathrm{d} y = \int_{\mathbb{R} } LG(x,y)g\mathrm{d} y = \int_{\mathbb{R} } \delta (x-y)g(y)\mathrm{d} y=g(x)
.\]
Here we must commute the differential operator with the integral using some theorem from real analysis that does not always apply. If $g$ is suffiiently nice, then it does.

We apply this to our case. Define $G(x,y)=\widetilde{G} (x-y)$, and define $\widetilde{G} (x)=\frac{m}{2} e^{-m \left| x \right| } $ when $m\neq 0$, and $\widetilde{G} (x)=-\frac{1}{2} \left| x \right| $ when $m=0$. This is a Green's function for $\Delta+m^2$. for example, we check this when $m=0$. We will choose a compactly supported test function $g$ such that we can do integration by parts and the boundary terms will vanish. Additionally, copactly supported functions vanish at $\pm \infty $. We have
\begin{align*}
	\int_{\mathbb{R} } \Delta (\widetilde{G} )\cdot f \mathrm{d} x &=\int_{\mathbb{R} } \widetilde{G} \cdot \Delta f\mathrm{d} x\\
								       &=\int_{\mathbb{R} } -\frac{1}{2} \left| x \right| \left( -\frac{\partial ^2f}{\partial x^2}  \right) \mathrm{d} x\\
								       &=-\frac{1}{2} \left( \int_{-\infty }^{0} -x\left(- \frac{\partial ^2f}{\partial x^2}\right)   \mathrm{d} x + \int_{0}^{\infty } x \left( -\frac{\partial ^2f}{\partial x^2}  \right) \mathrm{d} x \right) \\
								       &=-\frac{1}{2} \left( \int_{-\infty }^{0} x \frac{\partial ^2f}{\partial x^2} \mathrm{d} x + \int_{\infty }^{0} x \frac{\partial ^2f}{\partial x^2} \mathrm{d} x \right) \\
								       &=-\frac{1}{2} \left( -\int_{-\infty }^{0} \frac{\partial f}{\partial x}  \mathrm{d} x - \int_{\infty }^{0} \frac{\partial f}{\partial x} \mathrm{d} x \right) \\
								       &=\frac{1}{2} 2f(0)=f(0)
.\end{align*}
Therefore there exists a Green's function for $\Delta $, and $\pi $ is indeed a quasi-isomorphism.
\end{proof}
